\documentclass[11pt,a4paper]{article}
\usepackage[margin=2.4cm]{geometry}
\usepackage{amsmath,amssymb}
\usepackage{booktabs}
\usepackage{microtype}
\usepackage[colorlinks=true,linkcolor=blue,urlcolor=blue]{hyperref}

\newcommand{\bs}{\mathbf{s}}
\newcommand{\bu}{\mathbf{u}}
\newcommand{\bn}{\mathbf{n}}
\newcommand{\bd}{\mathbf{d}}
\newcommand{\bT}{\mathbf{T}}
\newcommand{\bE}{\mathbf{E}}

\title{Panorama-anchored ray tracing:\\ what the method is, why it works, and where it breaks}
\author{Notes for Robin}
\date{2026-08-01}

\begin{document}
\maketitle

\begin{abstract}
This is an honest account of the co-located transmitter idea, written because the
external criticism of it was correct but badly explained. The short version: the
objection is real, it kills one \emph{claim} and not the \emph{method}, and the
repair costs nothing because it reuses the identical computation. Separately, the
method is genuinely clever, but for a different reason than ``it avoids ray
tracing''. The reason it is clever is that the closed monostatic paths form a
measure-zero set, which makes them invisible to Monte Carlo ray shooting and
findable by a direct search. All numbers below come from this repository.
\end{abstract}

\section{What we have at every point}

Stand at a point $\bs$ where a street-level panorama was captured. Three
independent sources of information are already registered to each other there.

\begin{enumerate}
\item \textbf{The panorama.} At Korenmarkt the stitched Street View equirectangular
      image is $16384\times 8192$, which is $45.5$ pixels per degree.
\item \textbf{The semantic segmentation.} Every pixel carries a class, and through
      the concept catalogue a distribution over RF materials and attributes. The
      class decision grid is currently about $4.3$ pixels per degree, set by the
      segmentation network's mask resolution rather than by the image.
\item \textbf{The support mesh.} Photogrammetric city tiles, $157\,744$ triangles
      inside a $130$~m crop at Korenmarkt, with a stated geometric error of
      $2.006$~m.
\end{enumerate}

Write $\bu(\theta,\phi)$ for the unit direction at spherical angles
$(\theta,\phi)$. For each such direction we can read off, with no ray casting at
all:
\begin{equation}
  \bu \;\longmapsto\; \bigl(\, R(\bu),\; \bn(\bu),\; c(\bu),\; p_{\text{mat}}(\bu) \,\bigr),
\end{equation}
the first-hit range from the mesh depth buffer, the surface normal from the mesh,
the semantic class from the segmentation, and the material posterior from the
concept catalogue. The hit point itself is $\bs + R(\bu)\,\bu$.

The mesh and the image are good at complementary things, and it is worth being
precise about the asymmetry, because it is the real justification for the whole
approach.

\begin{table}[h]
\centering
\begin{tabular}{lrrr}
\toprule
Range & one degree subtends & one semantic cell & mesh geometric error \\
\midrule
$10$~m  & $0.17$~m & $4.1$~cm  & $200$~cm \\
$40$~m  & $0.70$~m & $16.2$~cm & $200$~cm \\
$100$~m & $1.75$~m & $40.6$~cm & $200$~cm \\
\bottomrule
\end{tabular}
\caption{Angular resolution of the semantic layer against the geometric error of
the support mesh. The image resolves lateral structure roughly an order of
magnitude finer than the mesh resolves position.}
\end{table}

More importantly, on the material axis the comparison is not a ratio. The mesh
carries \emph{no} material information whatsoever. Every facade exported from
tiles alone is an untextured triangle with no permittivity, no roughness, and no
thickness. The panorama is the only source of any of it, and it is also the only
source of clutter: poles, signs, bollards, railings and vehicles are absent from
the photogrammetry but visible in the image.

\section{The first bounce is free, and what that is actually worth}

A conventional shooting-and-bouncing tracer spends a bounding-volume-hierarchy
traversal, $O(\log M)$ in the triangle count $M$, to find each ray's first
intersection. Here the first intersection is a texture lookup, $O(1)$, because
the panorama \emph{is} the first-hit buffer.

It is tempting to sell this as the speedup, and I want to be blunt that it is
not much of one. For a three-bounce budget you replace one traversal out of
three, so the saving is a factor of roughly $1.5$, not orders of magnitude. If
someone claims this method is fast because the first bounce is free, they are
overselling it.

What the free first bounce is actually worth is \emph{data quality where the
physics is most sensitive}. Each additional bounce costs a reflection
coefficient. For brick at $28$~GHz, with $\varepsilon_r = 3.910 - j0.0260$
evaluated from ITU-R P.2040-4, the specular loss is

\begin{table}[h]
\centering
\begin{tabular}{lrr}
\toprule
Incidence $\theta_i$ & $|R_{\text{TE}}|$ & $|R_{\text{TM}}|$ \\
\midrule
$0^\circ$  & $0.328$ \;($-9.68$~dB)  & $0.328$ \;($-9.68$~dB) \\
$30^\circ$ & $0.377$ \;($-8.48$~dB)  & $0.278$ \;($-11.12$~dB) \\
$60^\circ$ & $0.561$ \;($-5.02$~dB)  & $0.048$ \;($-26.46$~dB) \\
$80^\circ$ & $0.816$ \;($-1.77$~dB)  & $0.433$ \;($-7.28$~dB) \\
\bottomrule
\end{tabular}
\caption{Fresnel magnitudes for brick at $28$~GHz, computed with this repository's
material library.}
\end{table}

So a second bounce is typically $5$ to $10$~dB down and a third is $10$ to
$20$~dB down. The first bounce carries most of the power, and the first bounce is
exactly the one where we have centimetre-scale material and clutter information
instead of a $2$~m-error untextured triangle. The design is well matched to the
physics: our best data sits on the term that dominates.

Note in passing the $60^\circ$ row. The two polarisations differ by $21$~dB near
the pseudo-Brewster angle. Any scalar treatment of this scene is throwing away a
first-order effect, which is the concrete argument for carrying a
polarisation-complete field rather than a scalar amplitude.

\section{Why this is a search and not a ray trace}

This is the part that is genuinely clever, and it deserves a careful statement
because it is easy to get backwards.

\subsection{One bounce}

A ray leaves $\bs$, strikes a surface at $A$, and returns to $\bs$. Since the
transmitter and receiver coincide, the outgoing and return segments are the same
segment traversed twice. This is pure backscatter, and it is not a search problem
at all: it is an integral over the panorama. Every pixel contributes, weighted by
its backscattering coefficient at its own incidence angle. No ray casting, no
visibility test, because $\bs$ sees $A$ by construction. The entire one-bounce
term is a weighted sum over the image.

\subsection{Two bounces, and the measure-zero argument}

Now consider $\bs \to A \to B \to \bs$ with specular reflection at both surfaces.
Fix the launch direction $\bu$. The hit point $A$ and its normal $\bn_A$ follow
from the panorama. The reflected direction is then determined:
\begin{equation}
  \bd_1 = \bu - 2(\bu\cdot\bn_A)\,\bn_A .
\end{equation}
Tracing along $\bd_1$ gives $B$ and its normal $\bn_B$. For the path to close,
the reflection at $B$ must send the ray precisely back to $\bs$:
\begin{equation}
  \bd_2 = \bd_1 - 2(\bd_1\cdot\bn_B)\,\bn_B
  \quad\text{must equal}\quad
  \frac{\bs - B}{\lVert \bs - B\rVert} .
  \label{eq:closure}
\end{equation}

Count the degrees of freedom. The launch direction $\bu$ supplies two free
parameters. Equation~\eqref{eq:closure} is an equality between two unit vectors,
which is two scalar constraints. Two parameters, two constraints: the solutions
are \emph{isolated points} in the launch sphere. Not an empty set, and not a
region. Isolated points.

That is the whole insight, and it has a sharp practical consequence.

\begin{quote}
A set of isolated points has zero measure on the sphere. A Monte Carlo ray
shooter samples the launch sphere at random and therefore finds a closed
two-bounce specular path with probability zero. It will never find one, no matter
how many rays you shoot.
\end{quote}

So for this geometry, ordinary ray shooting is not merely inefficient. It is
structurally incapable of returning the answer. Something else is required, and
the something else is a direct search.

\subsection{How the search actually works}

The classical construction is the image method, and it turns the search into an
$O(1)$ test per pair of surfaces. Given candidate surfaces $A$ and $B$:

\begin{enumerate}
\item Mirror $\bs$ in the plane of $A$ to get $\bs'$.
\item Mirror $\bs'$ in the plane of $B$ to get $\bs''$.
\item The straight segment from $\bs''$ to $\bs$ crosses the plane of $B$; that
      crossing is the candidate reflection point on $B$. Unfolding once more gives
      the candidate point on $A$.
\item Accept the path if both points lie inside their finite triangles, if both
      reflections are on the correct side, and if all three segments are
      unoccluded.
\end{enumerate}

Each ordered pair yields at most one candidate path, tested in constant time. So
the two-bounce problem is a loop over pairs of visible surface elements, which is
exactly the framing you arrived at independently, and the degree-of-freedom count
above is why it is the right framing rather than merely a convenient one.

This is also where the fishnet earns its place. The search is quadratic in the
number of surface elements, so element count is the cost driver:

\begin{center}
\begin{tabular}{lrr}
\toprule
Representation & elements, four crops & ordered pairs \\
\midrule
Image-space quadtree & $257\,310$ & $6.6\times 10^{10}$ \\
Fishnet cutter (measured) & $10\,534$ & $1.1\times 10^{8}$ \\
Visible support triangles (floor) & $5\,351$ & $2.9\times 10^{7}$ \\
\bottomrule
\end{tabular}
\end{center}

These are measured at Korenmarkt, not estimated. A factor of $24$ in element
count is a factor of $596$ in pair count, and it costs $0.19\%$ of source pixels
changing class. The argument for collapsing flat facades into few large polygons
is therefore not aesthetic and not about file size. It is what decides whether
the search is tractable.

Note the third row. A triangle lying entirely inside one semantic island is
preserved untouched, so the cutter cannot emit fewer faces than there are visible
support triangles. It currently sits a factor of $2$ above that floor, which means
the remaining headroom is not in the cutter at all: it is in the support mesh,
whose $2$~m-error triangulation is finer than the planar structure it represents.
Merging coplanar tile triangles before cutting is the next lever, and it is worth
roughly another factor of $4$ in pair count for free.

\subsection{Three bounces}

For $\bs \to A \to B \to C \to \bs$, both $A$ and $C$ are visible from $\bs$ and
so are read from the panorama, while $B$ need not be visible. The same image
construction applies with one more unfolding, and the degree-of-freedom count now
has four free parameters against two constraints, so solutions form a
two-parameter \emph{family} rather than isolated points. Three-bounce specular
returns are therefore generically abundant where two-bounce ones are rare.

There is a neater way to see the same thing. Reflect the outgoing direction at
$A$, and independently reflect at $C$ the direction from $\bs$ to $C$. Both
reflected rays must pass through $B$, so $B$ lies at their intersection, and the
required normal at $B$ follows from the two directions meeting there. Their
closest approach distance is a physical residual that measures how well the
hypothesis holds. External commentary treated $B$ as an unknown requiring a
learned scene-completion model. We do not need that. We have the support mesh, so
$B$ is a ray cast. What we genuinely lack at $B$ is not its geometry but its
\emph{material}, because $B$ is out of view and carries no image evidence. That is
a much smaller and much sharper problem than inventing a hidden city.

\subsection{Where this argument stops applying}

Everything in this section is conditional on one thing: that the path must
\emph{return to a point}. The closure condition \eqref{eq:closure} is an equality
between two unit vectors, and it is an equality only because the transmitter and
the receiver are the same point. That is what makes the solution set isolated, and
that is what makes random shooting useless.

Section~\ref{sec:repair} replaces that requirement. There the quantity of interest
is not the returning field but the field associated with rays that \emph{escape}
in direction $\Omega$, and $\Omega$ is accumulated into a finite angular bin. A
bin has non-zero solid angle. There is no closure constraint at all, only a
half-open one: leave $\bs$, interact, exit somewhere. So a randomly shot ray lands
in some bin with probability one, and the measure-zero obstruction simply is not
there.

The honest consequence is that the search is \emph{not} the primary estimator. The
primary estimator in the adjoint frame is ordinary gather-based Monte Carlo with
escaping rays recorded, which is a well-understood thing rather than a novel one.
The pair search survives in two narrower roles, both real:

\begin{itemize}
\item As the deterministic specular branch. Mirror lobes on smooth planes are
      narrow, and narrow lobes are exactly what Monte Carlo samples badly. Enumerating
      them in closed form over merged planes and treating the rest stochastically is
      the standard split, and it is where the image method belongs.
\item As the derived monostatic diagnostic, where the closure constraint genuinely
      does apply and where the measure-zero argument is therefore exactly right. If
      that number is ever quoted, the deterministic branch is the only way to obtain
      its specular part.
\end{itemize}

This also means the element-count table above should not be read as the argument
for the fishnet. Quadratic pair enumeration is a cost only in the specular branch.
The fishnet's load-bearing justifications are the ones that survive the change:
per-element material and roughness bookkeeping, the solid-angle pruning bound, and
keeping the number of things carrying a semantic posterior small enough to store
for ten cities. Those hold in either frame. The $596\times$ figure holds only in
the branch that enumerates pairs.

\section{What the search misses, and it is most of the power}

Everything above assumes mirror-like reflection. Real surfaces at millimetre
wavelengths are not mirrors, and the numbers are not marginal.

The Rayleigh roughness parameter compares surface height deviation to
wavelength. At $28$~GHz the free-space wavelength is $10.71$~mm. An earlier draft
of this section assumed $1$~mm of RMS height on brick, which gives, from this
repository's implementation:

\begin{table}[h]
\centering
\begin{tabular}{lrrr}
\toprule
Incidence & Rayleigh $g$ & scattering coeff.\ $S$ & specular power retained $1-S^2$ \\
\midrule
$0^\circ$  & $1.174$ & $0.865$ & $0.25$ \\
$30^\circ$ & $1.016$ & $0.803$ & $0.36$ \\
$60^\circ$ & $0.587$ & $0.540$ & $0.71$ \\
$80^\circ$ & $0.204$ & $0.202$ & $0.96$ \\
\bottomrule
\end{tabular}
\caption{An \emph{assumed} millimetre of roughness on brick at $28$~GHz. The
number in the first column was never sourced, and the measurement campaign
described below says it is wrong by a factor of thirty.}
\end{table}

That millimetre does not survive contact with the metrology. Every direct
measurement of a brick face returns $0.024$ to $0.095$~mm, which is a Rayleigh
$g$ near $0.03$ and leaves the face $99.9\%$ specular. The same holds across the
material set: measured RMS heights for glass, render, as-cast concrete, wood and
dressed stone all sit between $0.001$ and $0.3$~mm, and every one of them is
mostly specular at $28$~GHz. Full table and provenance in \texttt{ROUGHNESS.md}.

So the conclusion this section reached may still be right while its mechanism is
wrong, and that distinction matters more than the number. Measured $28$~GHz
campaigns do find roughly a fifth of received power in the diffuse component. But
micro-roughness cannot produce it: applying the Rayleigh closure to a study's own
profilometry under-predicts that same study's own fitted scattering coefficient by
five to twenty-four times in amplitude and by twenty-nine to five-hundred-fifty
times in power, consistently, across five samples measured both ways. Whatever
generates the diffuse power at these frequencies, it is not the surface finish.

The candidate that fits the magnitude is structure rather than roughness: mortar
joints, unit-to-unit offsets and block relief, at a pitch of seven to fifteen
wavelengths. That is a grating with on the order of fifteen propagating orders,
not a smooth lobe, and the two readings imply genuinely different architectures.
A comb and a lobe carrying equal total power are indistinguishable in any
measurement fitted with a lobe in the first place, which is how the question has
stayed open. Settling it needs a bistatic $28$~GHz scan of one square metre of
real brickwork at fine enough angular resolution to tell the two apart. Nobody
appears to have done it, and it is an afternoon of anechoic time.

Either way it settles the architecture, and it settles it against a pure search:

\begin{itemize}
\item \textbf{Specular contributions} live on isolated points. Monte Carlo cannot
      find them. The image-method search must be used.
\item \textbf{Diffuse contributions} live on a continuum with a finite-width lobe.
      They are found perfectly well by ordinary Monte Carlo with a gather step at
      each bounce, and a search would be the wrong tool.
\end{itemize}

So the honest answer to ``does this replace normal ray tracing'' is: no, it
replaces the part of it that never worked, and keeps the part that does. The
deliverable is a hybrid. That is a more defensible position than either extreme,
and it is worth stating plainly in the paper rather than being discovered by a
reviewer.

\section{Knowing when to stop}

Both the search and the Monte Carlo need a termination rule, and the semantic
layer supplies an unusually principled one.

For a ray subtending a fixed solid angle $\Omega$, an \emph{extended} surface such
as a facade has an illuminated footprint growing as $\Omega R^2$. Substituting
$\sigma = \sigma_0 \Omega R^2$ into the radar equation gives
\begin{equation}
  \frac{P_r}{P_t} \;\propto\; \frac{\lambda^2 \sigma_0 \Omega}{(4\pi)^3 R^{2}},
\end{equation}
so the return falls as $R^{-2}$ and a dynamic range budget $D$ in dB maps through
$20\log_{10}$. A \emph{discrete} scatterer whose cross-section does not grow with
range, such as a bollard, a lamp post or a person, keeps the textbook $R^{-4}$ and
maps through $40\log_{10}$. Same budget, two different cull distances:

\begin{table}[h]
\centering
\begin{tabular}{lrr}
\toprule
Dynamic range $D$ & facade, $R^{-2}$ & discrete clutter, $R^{-4}$ \\
\midrule
$20$~dB & $100$~m  & $32$~m  \\
$30$~dB & $316$~m  & $56$~m  \\
$40$~dB & $1000$~m & $100$~m \\
\bottomrule
\end{tabular}
\caption{Maximum useful range for a contribution referenced to $10$~m.}
\end{table}

The point is not the specific numbers. It is that the segmentation already knows
which regime each direction is in, so one threshold produces a per-class cull
distance instead of a single compromise. A twin without semantics cannot do this.
It is the first place in the design where the semantic layer buys propagation
accuracy rather than only material labels.

Two practical notes. First, the cull must be computed from an \emph{upper bound}
on what a direction could contribute, using the most reflective material still
admissible under that pixel's posterior and the shortest geometric path. Then a
discarded swath is provably negligible rather than probably negligible, and the
threshold can be defended instead of tuned. Second, reflection losses can be
folded into an effective range, $r_{\text{eff}} = r/\sqrt{\eta}$ for a power
factor $\eta$, so one budget prunes uniformly across bounce orders. That bookkeeping
is valid for pruning bounds only. It must never enter the field sum, because a
complex reflection coefficient carries phase that no real distance can represent.

\section{The objection, stated properly}

Here is the criticism you found unclear, in what I hope is plainer form.

The environment is a linear system. Write $G(\mathbf{r}_1, \mathbf{r}_2)$ for the
field produced at $\mathbf{r}_1$ by a unit source at $\mathbf{r}_2$. A realistic
exposure question asks for $G(\bs, \mathbf{x})$, the field at the pedestrian
$\bs$ due to a base station at some other place $\mathbf{x}$. The co-located
construction computes $G(\bs,\bs)$, the response at $\bs$ to a source at $\bs$.
That is one entry of the operator, the diagonal one.

The concrete reason the diagonal does not give you the rest: consider the
three-bounce loop $\bs \to A \to B \to C \to \bs$. Its coefficient factorises,
roughly, as
\begin{equation}
  h_{\text{loop}} \;\approx\; \underbrace{h_{\bs \leftarrow B}^{(C)}}_{\text{return leg}}
  \;\cdot\; \Gamma_B \;\cdot\;
  \underbrace{h_{B \leftarrow \bs}^{(A)}}_{\text{outbound leg}} .
\end{equation}
Reciprocity tells you the two legs are equal, so what you measure is essentially
one leg \emph{squared}. But exposure from a real base station needs one leg
\emph{alone}, from the mast to $\bs$, and knowing a product does not determine
its factors when many different pairs give the same product. A strong loop can
arise from a strong outbound leg or from a strong return leg, and the co-located
measurement cannot tell you which.

Two concrete cases where the ranking inverts:

\begin{description}
\item[Open street, high exposure, low loop.] A pedestrian on a wide straight
      street with clear line of sight to a rooftop mast $80$~m away. Exposure is
      high and dominated by the direct term. But the geometry sends almost nothing
      back to the pedestrian's own position, so the monostatic return is small.
\item[Enclosed courtyard, low exposure, high loop.] A pedestrian in a small
      courtyard surrounded by close walls, with no mast in view. Multiple surfaces
      return energy strongly to the centre, so the monostatic return is large. But
      no external radiation reaches the courtyard, so the exposure is low.
\end{description}

The monostatic descriptor ranks these two backwards. It is therefore not an upper
bound on exposure and not a worst case, because it is not conservative in either
direction. It is a different quantity: a measure of how reverberant a location is.
That is a legitimate and interesting thing to compute. It simply is not exposure,
and calling it a representative worst case invites exactly the rebuttal above.

\section{The repair, which costs nothing}
\label{sec:repair}

Here is why none of this threatens the actual work.

Currently the trace shoots rays from $\bs$ and keeps the ones that come back.
Almost none come back, so almost every ray is discarded. Instead, keep every ray
that \emph{escapes} the local neighbourhood after up to three interactions, and
record its exit direction. Lorentz reciprocity says the reverse of each escaping
path is a valid incoming path from an external source. What you build is
\begin{equation}
  \bT_{\bs}(\Omega, \tau, f) \;=\;
  \text{field at } \bs \text{ from a unit plane wave arriving from } \Omega
  \text{ at delay } \tau,
\end{equation}
and then, for any deployment described by an angular incident power distribution
$\mathbf{Q}_{\bs}$,
\begin{equation}
  \mathbb{E}\bigl[\lVert \bE(\bs,f)\rVert^2\bigr]
  \;=\;
  \int_{4\pi}
  \operatorname{Tr}\!\left[
    \bT_{\bs}\,\mathbf{Q}_{\bs}\,\bT_{\bs}^{H}
  \right] d\Omega .
  \label{eq:factorise}
\end{equation}

Same rays. Same free first hit. Same cull. Same materials. The only change is
which rays you write down at the end. The monostatic response is still available,
as the subset of paths that happen to close, and becomes a derived diagnostic
rather than the headline.

Three things follow that are worth more than the loop was.

\paragraph{The antenna zoo problem disappears.} $\bT_{\bs}$ never mentions a
transmitter. It is the environment's response to a plane wave. Base stations enter
only through $\mathbf{Q}_{\bs}$, which is an angular \emph{distribution} carrying
density, EIRP, load and line-of-sight probability, not a catalogue of hardware.
You never have to describe a 6G antenna that does not exist yet. You describe
where power arrives from.

\paragraph{Reproducibility, which was the original complaint.} Equation
\eqref{eq:factorise} separates an expensive, deployment-independent, per-location
quantity from a cheap, assumption-laden one. Publish $\bT_{\bs}$ for ten cities
once. A reviewer who dislikes your deployment substitutes their own
$\mathbf{Q}_{\bs}$ and redoes the integral in a notebook, without re-tracing
anything.

\paragraph{A defensible worst case actually exists.} Take
\begin{equation}
  \chi_{\max}(\bs) \;=\; \max_{\Omega}\;
  \sigma_{\max}^2\bigl(\bT_{\bs}(\Omega)\bigr),
\end{equation}
the largest singular value over all arrival directions. This answers ``from the
worst possible external direction, how much does this environment enhance the
incident field here''. It is a genuine upper bound over deployments at fixed
power, it is one scalar per location, and the maximum over $\Omega$ dominates any
$\mathbf{Q}_{\bs}$ a reviewer proposes. Alongside it, $\chi_{\text{iso}}$ gives
the isotropically illuminated average and $\chi_{\text{roof}}$ the rooftop-weighted
realistic case.

Line of sight also stops being bolted on. It is the zero-bounce term
$\bT_{\bs} = \bT_0 + \bT_{\text{MP}}$, and its angular support is exactly the set
of directions where the mesh first-hit range exceeds the source range, which the
depth buffer already gives you. The sky mask is the special case of infinite
first-hit range, which means line-of-sight probability is \emph{measured} at each
location rather than taken from a fitted model.

The one honest cost of the change: you must define the exit surface and confront
whether a mast at $50$ to $200$~m is far enough away for a clean plane-wave
treatment. The closed loop sidesteps that question. I think it is worth paying.

\section{Verdict}

Do I believe the objection? Yes, and it is not subtle once stated in terms of the
operator diagonal. But it invalidates the interpretation, not the machinery, and
the machinery is the part that took work.

Is the method clever? Yes, but not for the reason it is usually pitched, and not
for the reason I gave two drafts ago either. The pitch ``the first bounce is free
so it is fast'' is weak, worth maybe a factor of $1.5$. The pitch ``closed
specular paths are measure zero so pair search is forced'' is sharper, but it is
an argument about the co-located geometry specifically, and once the observable
moves to the adjoint tensor that geometry is gone and the argument goes with it.

What is left is the part that was load-bearing all along. The panorama supplies
material, roughness and clutter at centimetre scale on the bounce that carries
most of the power, where a mesh-only tracer has a $2$~m-error triangle and no
material at all. That advantage is indifferent to which observable you compute,
and it is the one thing here that a conventional tracer with the same mesh cannot
have.

Where I would temper the enthusiasm: measured $28$~GHz campaigns put roughly a
fifth of received power in the diffuse component, so the specular branch is a
branch of the estimator rather than the whole estimator, and the result is a
hybrid that the paper should describe as one. The larger honesty is that we do
not yet know what produces that diffuse fifth. It is not surface finish, because
every measured facade material is nearly specular at this wavelength. The
leading candidate is masonry structure acting as a coarse grating, and until
someone runs the bistatic scan described above, any scattering lobe we ship is a
fitted stand-in for a mechanism we have not identified.

\end{document}
